\documentclass[letterpaper]{article} 
\usepackage[utf8]{inputenc}
\linespread{0.85}
\usepackage[T1]{fontenc}
\usepackage{amsmath}
\usepackage{amsfonts}
\usepackage{amssymb}
\usepackage{array}
\usepackage{fontspec}
\usepackage{booktabs}
\usepackage{hyperref}
\usepackage[version=4]{mhchem}
\usepackage{stmaryrd}
\usepackage[dvipsnames]{xcolor}
\colorlet{LightRubineRed}{RubineRed!70}
\colorlet{Mycolor1}{green!10!orange}
\definecolor{Mycolor2}{HTML}{00F9DE}
\usepackage{graphicx}
\usepackage{amsmath}
\usepackage{graphicx}
\usepackage{capt-of}
\usepackage{lipsum}
\usepackage{fancyvrb}
\usepackage{tabularx}
\usepackage{listings}
\usepackage[export]{adjustbox}
\graphicspath{ {./images/} }
\usepackage[utf8]{inputenc}
\usepackage[english]{babel}
\usepackage{float}
\usepackage{lipsum}
\usepackage{graphicx}
\usepackage{float}
\usepackage[margin=0.7in]{geometry}
\usepackage{amsmath}
\usepackage{graphicx}
\usepackage{capt-of}
\usepackage{tcolorbox}
\usepackage{lipsum}
\usepackage{graphicx}
\usepackage{float}
\usepackage{listings}
\definecolor{deepred}{RGB}{139,0,0}
\usepackage{hyperref} 
\usepackage{xcolor} % For custom colors
\lstset{
	language=Python,                % Choose the language (e.g., Python, C, R)
	basicstyle=\ttfamily\small, % Font size and type
	keywordstyle=\color{blue},  % Keywords color
	commentstyle=\color{gray},  % Comments color
	stringstyle=\color{red},    % String color
	numbers=left,               % Line numbers
	numberstyle=\tiny\color{gray}, % Line number style
	stepnumber=1,               % Numbering step
	breaklines=true,            % Auto line break
	backgroundcolor=\color{black!5}, % Light gray background
	frame=single,               % Frame around the code
}
\usepackage{float}
\usepackage[]{amsthm} %lets us use \begin{proof}
\usepackage[]{amssymb} %gives us the character \varnothing

	\title{Homework 2, MATH 4155}
	\author{Zongyi Liu}
	\date{Sun, Feb 15, 2026}
	\begin{document}
		\maketitle
		
		\section{Question 1}
		
		Construct a distribution function which is discontinuous at every rational point, and continuous at all irrational points on the real line. Conversely: is there a distribution function which is discontinuous at every irrational point, and continuous at all rational points on the real line?
		
		\textbf{Answer}
		
		\clearpage
		
	
	\section{Question 2}
	
	{\color{deepred}\fontspec{Georgia} Distribution and CDF}
	
	\begin{enumerate}
		
		
		
		\item[(i)] Suppose $\mu$ is a probability measure on the Borel subsets of the real line.
		We use it to define a function $F : \mathbb{R} \to [0,1]$ via
		\[
		F(x) := \mu\bigl((-\infty,x]\bigr), \qquad x \in \mathbb{R}.
		\]
		Show that this function is nondecreasing and right continuous; that it satisfies
		$F(-\infty)=0$ and $F(\infty)=1$; and that for any real numbers $a<b$ we have
		\begin{equation}
			\mu((a,b]) = F(b)-F(a), \qquad
			\mu([a,b)) = F(b-)-F(a-), \qquad
			\mu(\{a\}) = F(a)-F(a-),
			\tag{0.1}
		\end{equation}
		and
		\begin{equation}
			\mu([a,b]) = F(b)-F(a-), \qquad
			\mu((a,b)) = F(b-)-F(a).
			\tag{0.2}
		\end{equation}
		
		\item[(ii)] Given a random variable $X:\Omega \to \mathbb{R}$ on a probability space
		$(\Omega,\mathcal{F},\mathbb{P})$, consider the induced measure
		$\mu_X = \mathbb{P}\circ X^{-1}$.
		Show that the function
		\[
		F_X(x) := \mu_X\bigl((-\infty,x]\bigr) = \mathbb{P}(X \le x),
		\qquad x \in \mathbb{R},
		\]
		is nondecreasing and right continuous, and satisfies
		$F_X(-\infty)=0$, $F_X(\infty)=1$ as well as the properties
		\textup{(0.1)}, \textup{(0.2)} above.
	\end{enumerate}
		
		\textbf{Answer}
		
		\clearpage
		
		\section{Question 3}
		
		{\color{deepred}\fontspec{Georgia} De Moivre-Laplace for Coin Tossing} 
		
		Let $X_1,X_2,\dots$ be independent random variables with
			\[
			\mathbb{P}(X_j=1)=1-\mathbb{P}(X_j=0)=p\in(0,1)
			\quad \text{for all } j\in\mathbb{N},
			\]
			and denote by
			\[
			S_n=\sum_{j=1}^n X_j
			\]
			the ``number of successes in the first $n$ tosses.'' Show that
			\begin{equation}
				\mathbb{P}\!\left[
				a \le \frac{S_n-np}{\sqrt{np(1-p)}} \le b
				\right]
				\;\longrightarrow\;
				\Phi(b)-\Phi(a),
				\qquad a<b \text{ in } \mathbb{R},
				\tag{0.3}
			\end{equation}
			as $n\to\infty$, where
			\[
			\Phi(x):=\int_{-\infty}^x \varphi(\xi)\,d\xi,
			\qquad
			\varphi(x):=\frac{e^{-x^2/2}}{\sqrt{2\pi}},
			\]
			is the so-called \emph{standard Normal distribution function}.
			\medskip
			
			\noindent
			(\emph{Hint:} With the help of the Stirling formula
			\[
			n! \sim \sqrt{2\pi n}\,n^n e^{-n},
			\]
			actually in its stronger form
			\begin{equation}
				n! = \sqrt{2\pi}\, n^{\,n+\frac12} e^{-n+\varepsilon_n}
				\qquad \text{with} \qquad
				\frac{1}{12n+1}<\varepsilon_n<\frac{1}{12n},
				\tag{0.4}
			\end{equation}
			from the previous assignment, show first the ``local'' form
			\[
			\lim_{n\to\infty}
			\left(
			\sqrt{npq}\,
			\frac{n!}{k_n!(n-k_n)!}\,
			p^{k_n}(1-p)^{n-k_n}
			\right)
			=
			\frac{e^{-x^2/2}}{\sqrt{2\pi}}
			=\varphi(x),
			\]
			of this result, where $x\in\mathbb{R}$ is fixed and
			\[
			k_n := x\sqrt{np(1-p)}+np;
			\]
			then observe that this convergence is uniform over $x$ in the bounded
			interval $[a,b]$, as the next step towards \textup{(0.3)}.)
		
		
		\textbf{Answer}
		
		\clearpage
		
		\section{Question 4}
		Let $(\Omega, \mathcal{F}, \mathbb{P})$ be a probability space,
			and consider a set $\mathcal{J}$ which is (at most) countable.
			For every $j \in \mathcal{J}$ we are given two events
			$B_j \subseteq A_j$ in $\mathcal{F}$. Show that
			\[
			\mathbb{P}\!\left(\bigcup_{j \in \mathcal{J}} A_j\right)
			-
			\mathbb{P}\!\left(\bigcup_{j \in \mathcal{J}} B_j\right)
			\;\le\;
			\sum_{j \in \mathcal{J}}
			\bigl[\mathbb{P}(A_j) - \mathbb{P}(B_j)\bigr].
			\]
			
		
		\textbf{Answer}
		
		\clearpage
		
		\section{Question 5}
		{\color{deepred}\fontspec{Georgia} Coin Tossing} 
		
		Use (0.3) to establish the \textsc{Bernoulli Weak Law of Large Numbers}
		\[
		\lim_{n \to \infty}
		\mathbb{P}\!\left(\left|\overline{X}_n - p\right| > \varepsilon\right)
		=
		0,
		\qquad \forall \,\varepsilon > 0,
		\]
		in the Coin-Tossing context of Exercise 3.
		Here $\overline{X}_n := S_n/n$ is the relative frequency of $1$’s
		in the first $n$ independent tosses of the coin.
		
		
		\textbf{Answer}
		
		\clearpage
		
		\section{Question 6}
		
		{\color{deepred}\fontspec{Georgia} Walsh 3.46} 
		
		Let $X_1, X_2, \ldots, X_n$ be independent random variables with mean zero and
		variance one. Let $Y$ be a square-integrable random variable.
		Let $a_j = E\{Y X_j\}$.
		Prove Bessel's inequality:
		\[
		E\{Y^2\} \ge \sum_{j=1}^n a_j^2.
		\]
		Show that Bessel's inequality holds if the $X_j$ are only uncorrelated,
		not necessarily independent.
		Show that this extends to an infinite sequence of $X_n$.
		
		\medskip
		
		\textit{Hint:} Take the expectation of
		\[
		\left( Y - \sum_{j=1}^n E\{Y X_j\} X_j \right)^2.
		\]
		
		
		
		\textbf{Answer}
		
		\clearpage
		
		
		\section{Question 7}
		
		
		{\color{deepred}\fontspec{Georgia} Walsh 3.47} 
		
		Let $X_1, X_2, \ldots$ be independent mean zero, variance one random variables.
		Let $Y$ be a square-integrable random variable.
		Show that
		\[
		E\{Y X_n\} \longrightarrow 0
		\quad \text{as } n \to \infty.
		\]
		
		\medskip
		
		\textit{Hint:} Use the previous problem.
		
		
		\textbf{Answer}
		
		\clearpage
		
		\section{Question 8}
		
		
		{\color{deepred}\fontspec{Georgia} Walsh 3.37}
		 
		Two light fixtures hang in different hallways of an apartment building.
		For each fixture, a given light bulb lasts an exponential time with mean
		$\lambda > 0$ before it has to be replaced ($\lambda$ is the same for both).
		There is a stack of $n$ replacement bulbs in a closet in hallway $A$,
		and $k$ replacement bulbs in the corresponding closet in hallway $B$.
		What is the probability that hallway $A$ runs out of bulbs before $B$ does?
		
		\textbf{Answer}
		
		\clearpage
		
		\section{Question 9}
		
		{\color{deepred}\fontspec{Georgia} Walsh 2.51} 
		
		It is said that many are called and few are chosen. Suppose that the number called is Poisson $(\lambda)$. Each person called is chosen, independently, with probability $p$. Show that the distribution of the number chosen is Poisson $(p\lambda)$.
		
		
		\textbf{Answer}
		
		\clearpage
		
		\section{Question 10}
		
		{\color{deepred}\fontspec{Georgia} Binomial Probability} 
		
		In the context of Exercise 3, compute the probability
		\[
		\mathbb{P}(S_n \text{ is even}).
		\]
		(\textit{Hint}: Recall the Binomial Theorem $(x + y)^n = \sum_{k=0}^n (n!/k!(n - k)!) \, x^k y^{n-k}$.)
		
		
		\textbf{Answer}
		
		\clearpage
		
		
		
		
	\end{document}
